% Options for packages loaded elsewhere
\PassOptionsToPackage{unicode}{hyperref}
\PassOptionsToPackage{hyphens}{url}
\PassOptionsToPackage{dvipsnames,svgnames,x11names}{xcolor}
\documentclass[
  a4paper,
]{article}
\usepackage{xcolor}
\usepackage[margin=25mm]{geometry}
\usepackage{amsmath,amssymb}
\setcounter{secnumdepth}{-\maxdimen} % remove section numbering
\usepackage{iftex}
\ifPDFTeX
  \usepackage[T1]{fontenc}
  \usepackage[utf8]{inputenc}
  \usepackage{textcomp} % provide euro and other symbols
\else % if luatex or xetex
  \usepackage{unicode-math} % this also loads fontspec
  \defaultfontfeatures{Scale=MatchLowercase}
  \defaultfontfeatures[\rmfamily]{Ligatures=TeX,Scale=1}
\fi
\usepackage{lmodern}
\ifPDFTeX\else
  % xetex/luatex font selection
  \setmainfont[]{DejaVu Serif}
  \setsansfont[]{DejaVu Sans}
  \setmonofont[]{DejaVu Sans Mono}
\fi
% Use upquote if available, for straight quotes in verbatim environments
\IfFileExists{upquote.sty}{\usepackage{upquote}}{}
\IfFileExists{microtype.sty}{% use microtype if available
  \usepackage[]{microtype}
  \UseMicrotypeSet[protrusion]{basicmath} % disable protrusion for tt fonts
}{}
\makeatletter
\@ifundefined{KOMAClassName}{% if non-KOMA class
  \IfFileExists{parskip.sty}{%
    \usepackage{parskip}
  }{% else
    \setlength{\parindent}{0pt}
    \setlength{\parskip}{6pt plus 2pt minus 1pt}}
}{% if KOMA class
  \KOMAoptions{parskip=half}}
\makeatother
\usepackage{longtable,booktabs,array}
\newcounter{none} % for unnumbered tables
\usepackage{calc} % for calculating minipage widths
% Correct order of tables after \paragraph or \subparagraph
\usepackage{etoolbox}
\makeatletter
\patchcmd\longtable{\par}{\if@noskipsec\mbox{}\fi\par}{}{}
\makeatother
% Allow footnotes in longtable head/foot
\IfFileExists{footnotehyper.sty}{\usepackage{footnotehyper}}{\usepackage{footnote}}
\makesavenoteenv{longtable}
\setlength{\emergencystretch}{3em} % prevent overfull lines
\providecommand{\tightlist}{%
  \setlength{\itemsep}{0pt}\setlength{\parskip}{0pt}}
\usepackage{luatexja}
\usepackage{xurl}
\usepackage{bookmark}
\IfFileExists{xurl.sty}{\usepackage{xurl}}{} % add URL line breaks if available
\urlstyle{same}
\hypersetup{
  pdftitle={A Catalogue of Symmetries and Magic Numbers for Reading Out Spacetime from a State Space},
  pdfauthor={Noriaki Kihara (WF System Co., Ltd.) ORCID 0009-0004-6753-4020},
  colorlinks=true,
  linkcolor={Maroon},
  filecolor={Maroon},
  citecolor={Blue},
  urlcolor={Blue},
  pdfcreator={LaTeX via pandoc}}

\title{A Catalogue of Symmetries and Magic Numbers for Reading Out
Spacetime from a State Space}
\usepackage{etoolbox}
\makeatletter
\providecommand{\subtitle}[1]{% add subtitle to \maketitle
  \apptocmd{\@title}{\par {\large #1 \par}}{}{}
}
\makeatother
\subtitle{― The Answer-Check Standard of a Research Programme Taking
Anonymity and Self-Consistency as Axioms ―}
\author{Noriaki Kihara (WF System Co., Ltd.) ORCID 0009-0004-6753-4020}
\date{October 18, 2026 Version DOI 10.5281/zenodo.23444452}

\ltjsetparameter{jacharrange={-2,-3}}
\begin{document}
\maketitle

\textbf{Author:} Noriaki Kihara\\
\textbf{ORCID:} 0009-0004-6753-4020\\
\textbf{Version DOI:} 10.5281/zenodo.23444452\\
\textbf{Concept DOI:} 10.5281/zenodo.23444451\\
\textbf{Version:} v1.0\\
\textbf{Date:} 2026-10-18

\begin{center}\rule{0.5\linewidth}{0.5pt}\end{center}

\subsection{Abstract}\label{abstract}

The author's research programme asks whether time, space, mass, charge
and interactions can be read out from the relations of a set of complex
numbers \{zₙ\} alone, without placing a background spacetime, a ruler or
particle species from outside. The programme has published 156 papers so
far, but the numbering of axioms, the symbols and the claims differ
slightly from paper to paper, and it had become unclear against what
standard a given numerical experiment or thought experiment is to be
judged correct.

This paper presents no new derivation. Instead, it fixes a catalogue of
the symmetries the programme must satisfy and of the dimensionless
numbers (magic numbers) attached to those symmetries, giving every entry
a status (established, theory, candidate, open) and a provenance (a
published paper or external literature). The catalogue is a reference
document that determines which papers later papers should use as their
answer-check standard; that is why it contains many self-citations and
external citations.

The skeleton of the catalogue is as follows. (1) The zeroth axiom is
anonymity and the first axiom is self-consistency. As an axiom derived
from anonymity, we place the distinguishability axiom: the readout
cannot tell which layer a value belongs to. The squared quantity z′ = z²
is also just a complex number, so z and z′ cannot be distinguished in
the readout. On that basis the closure is defined by a single
definition, the linear zero closure of complex numbers Σzₙ = 0. The
square-sum closure Σzₙ² = Σz′ₙ = 0 is the same definition read at the z′
layer, not a new definition. The return of time evolution Uⁿ = I is a
theorem. (2) Below these we place the big-picture symmetries (time
symmetry of emission and absorption, local and global, the scale of
observation, electric--magnetic duality, Galois symmetry, conservation
of information), conjugacy and duality, and the projection symmetry. (3)
We set out a candidate axiom: an essential symmetry carries a
dimensionless magic number rather than the conserved quantity assumed by
Noether's theorem. (4) The requirements of the Standard Model, general
relativity and E₈, the algebra and geometry of this system, and the
locking symmetries of complex systems are listed in the same format. (5)
Magic numbers are divided into four systems: integer,
coupling-convention, irrational and derived.

We also present two mathematically confirmed results. A linear map that
preserves both the sum of true squares and the norm is necessarily a
real orthogonal map. The 24 units of the Hurwitz lattice D₄ split into
one vector set and two spinor sets (8 + 8 + 8), and triality cycles them
as the quaternion operation q ↦ a q a (a = (−1 + i + j + k)/2) while
preserving the 240 roots of E₈.

All 156 published papers are mapped to the rows of the catalogue: 81
rows have a published standard and 15 rows are set up for the first time
in this paper. The axiom numbering is reset in this paper, with a
correspondence table to the old numbering. The earlier claim that filled
the gap between 137 and α⁻¹ with a closed identity is revised to the
position that the gap itself is the content to be read. Numerical
disagreements, such as the rotation orders of the four-dimensional
lattice, the decomposition of 137 and the number of components of one
generation, are kept as open problems, with ``several interpretations
are possible'' as the present answer. Finally, we place the root
difficulty of the map from the state space to spacetime in the fact that
what is visible depends on the projection method, and we show that,
since spacetime is stably readable, the map must be at least a partial
projection.

\textbf{Keywords:} anonymity, self-consistency, zero square sum, linear
zero closure, distinguishability, finite order, symmetry, magic number,
Noether's theorem, conjugacy, duality, partial projection, triality, E₈,
mode locking

\begin{center}\rule{0.5\linewidth}{0.5pt}\end{center}

\subsection{1. Position of This Paper}\label{position-of-this-paper}

\subsubsection{1.1 Why a catalogue is
needed}\label{why-a-catalogue-is-needed}

The research programme published 156 papers between April and October
2026. Among the papers that state an axiom system alone, the Anonymous
Equal-Amplitude Composite-Wave Model: Basic Axiom System v9.1 {[}S1{]},
the symmetry generation from zero closure, finite order and
self-consistent geometry {[}S3{]}, the research project design {[}S14{]}
and the first and second thought experiments {[}S15, S16{]} each number
their axioms differently (§11). On the side of numerical experiments, no
document fixed the criterion for judging whether a result was derived
from the skeleton or came from parts imported from outside the skeleton.
In the author's own numerical experiments, a failure occurred in which
parts were swapped each time a result did not match until the experiment
``passed'' (§13.4).

This paper fixes that criterion. Whereas the research project design
{[}S14{]} fixes the direction of the project (goal image, prohibitions,
success conditions), this paper fixes the symmetries and numbers that
individual papers and numerical experiments must satisfy. The research
project design adopted a policy of no self-citation; this paper, on the
contrary, covers self-citations and external citations exhaustively,
because it determines which papers are to be used as the answer-check
standard.

\subsubsection{1.2 How statuses are
assigned}\label{how-statuses-are-assigned}

Each row of the catalogue is given one of the following four statuses.

{\def\LTcaptype{none} % do not increment counter
\begin{longtable}[]{@{}
  >{\raggedright\arraybackslash}p{(\linewidth - 2\tabcolsep) * \real{0.5000}}
  >{\raggedright\arraybackslash}p{(\linewidth - 2\tabcolsep) * \real{0.5000}}@{}}
\toprule\noalign{}
\begin{minipage}[b]{\linewidth}\raggedright
Status
\end{minipage} & \begin{minipage}[b]{\linewidth}\raggedright
Meaning
\end{minipage} \\
\midrule\noalign{}
\endhead
\bottomrule\noalign{}
\endlastfoot
Established & An established mathematical or physical fact, supported by
external literature or by the checks in this paper (§10) \\
Theory & A proposition set up within the theory of the programme, with a
published paper as provenance \\
Candidate & A conjectured correspondence, not yet checked \\
Open & Something that stands only as a question, including items
admitting several interpretations \\
\end{longtable}
}

\subsubsection{1.3 Resetting the
numbering}\label{resetting-the-numbering}

The numbering of this paper (A1--A6, A1′, Z1, W1--W6, K1--K7, P1,
S1--S20, E1--E6, G1--G18, C1--C4, L1--L9) is the standard for later
papers. The numbering of published papers is kept as numbering within
each paper and is not rewritten. The correspondence with the old
numbering is given in §11.1.

\begin{center}\rule{0.5\linewidth}{0.5pt}\end{center}

\subsection{2. Starting Point: The State Is Only a Set of Complex
Numbers}\label{starting-point-the-state-is-only-a-set-of-complex-numbers}

The state of this system is only a set of complex numbers

\[
\{z_n\},\qquad z_n=a_n+i\,b_n .
\]

In a sense it is a Hilbert space, and there are no positions, lattice
points, values sitting on vertices, fields or background spacetime. Each
zₙ is merely a point in the complex plane, and its magnitude

\[
R^2=a^2+b^2
\]

is read not as the radius of some space but as the magnitude of the
complex number itself.

Spacetime \{xyzt\}, distance, time and R are all read out from the
relations among the zₙ (ratios, phase differences, true squares, cross
terms). The largest misreading the author repeated during the research
was to believe that the background is xyz space. Under that misreading
one regards lattices and spheres as containers and tries to arrange
complex numbers in them. In this paper, lattices, S³ and Clifford
translations are treated as structures in phase space, not as
arrangements in xyz space.

The destination of the readout is split into the eight directions of two
quaternions:

\[
A=t+xi+yj+zk\quad(\text{readable side}),\qquad B=R+Q_1i+Q_2j+Q_3k\quad(\text{unreadable interior}) .
\]

The dimensions of number systems that preserve magnitude under squaring
are limited to 1, 2, 4 and 8 (Hurwitz's theorem), and eight directions
are the upper limit. Setting R and Q₁Q₂Q₃ up as independent axes avoids
the confusion by which Weyl's 1918 scale gauge theory made R (scale)
also carry the role of Q (charge) and reached an impasse. t is the clock
of space and R is the internal clock (mass), so two clocks can be read
separately.

\begin{center}\rule{0.5\linewidth}{0.5pt}\end{center}

\subsection{3. Axioms and Theorems}\label{axioms-and-theorems}

\subsubsection{3.0 The distinguishability axiom and the single
definition of
closure}\label{the-distinguishability-axiom-and-the-single-definition-of-closure}

{\def\LTcaptype{none} % do not increment counter
\begin{longtable}[]{@{}
  >{\raggedright\arraybackslash}p{(\linewidth - 4\tabcolsep) * \real{0.3333}}
  >{\raggedright\arraybackslash}p{(\linewidth - 4\tabcolsep) * \real{0.3333}}
  >{\raggedright\arraybackslash}p{(\linewidth - 4\tabcolsep) * \real{0.3333}}@{}}
\toprule\noalign{}
\begin{minipage}[b]{\linewidth}\raggedright
No.
\end{minipage} & \begin{minipage}[b]{\linewidth}\raggedright
Content
\end{minipage} & \begin{minipage}[b]{\linewidth}\raggedright
Kind
\end{minipage} \\
\midrule\noalign{}
\endhead
\bottomrule\noalign{}
\endlastfoot
A1′ & \textbf{Distinguishability axiom}: the readout cannot tell which
layer a read-out value belongs to. This applies anonymity (A1) to layers
rather than to elements & Axiom (derived from A1) \\
― & \textbf{Its appearance in the readout}: the squared quantity z′ = z²
is also just a complex number, so z and z′ cannot be distinguished in
the readout & Symmetry \\
― & \textbf{Choice of basic quantity}: the squared quantity z′ₙ = zₙ² is
also a basic quantity. Writing z′ = A + iB, A = a² − b² and B = 2ab &
Definition \\
Z1 & \textbf{Linear zero closure}: Σzₙ = 0 (Σaₙ = 0 and Σbₙ = 0). The
single definition of closure in this paper & Definition \\
― & \textbf{Its consequence}: the square-sum closure Σzₙ² = Σz′ₙ = 0 is
Z1 read at the z′ layer, not a new definition & Consequence \\
\end{longtable}
}

Under the distinguishability axiom, the one definition ``the sum of the
readable values is 0'' is imposed so that it holds whichever layer it is
read at. As a result, both Σzₙ = 0 and Σzₙ² = 0 follow. What looked like
two conditions is one definition seen in two readings. Together they
give

\[
\sum_n a_n=\sum_n b_n=0,\qquad \sum_n a_n^2=\sum_n b_n^2,\qquad \sum_n a_nb_n=0 ,
\]

and (Σaₙ)(Σbₙ) = 0 follows trivially. If the points (aₙ, bₙ) are placed
in the complex plane, their centroid is at the origin, their spreads
along the two axes are equal, and the two axes are uncorrelated. Up to
second order, the point set has no special origin and no special
direction, and no meaning is left in which axis is called real and which
imaginary. This is the first- and second-order expression, on the
complex plane, of Axiom 0 of the basic axiom system (no privileged axis,
sign or type) {[}S1{]}.

The linear sum is the lowest-degree quantity among those unchanged by
every permutation, so the definition of closure can be written with only
the lowest-degree invariant of anonymity and the indistinguishability of
layers. The squaring map z ↦ z² is not part of the definition; it stands
outside the definition as the map connecting the two readings.

Provenance: a precedent that used the linear closure as a condition on
parent states (centroid at the origin) is the complex-simplex study of
2026-08-30. The view that places the linear closure as the definition
and identifies it with the square-sum closure, and the
distinguishability axiom, come from the thought-experiment dialogue
record on which this paper is based (Appendix D).

\subsubsection{3.1 Two axioms}\label{two-axioms}

{\def\LTcaptype{none} % do not increment counter
\begin{longtable}[]{@{}
  >{\raggedright\arraybackslash}p{(\linewidth - 6\tabcolsep) * \real{0.2500}}
  >{\raggedright\arraybackslash}p{(\linewidth - 6\tabcolsep) * \real{0.2500}}
  >{\raggedright\arraybackslash}p{(\linewidth - 6\tabcolsep) * \real{0.2500}}
  >{\raggedright\arraybackslash}p{(\linewidth - 6\tabcolsep) * \real{0.2500}}@{}}
\toprule\noalign{}
\begin{minipage}[b]{\linewidth}\raggedright
No.
\end{minipage} & \begin{minipage}[b]{\linewidth}\raggedright
Axiom
\end{minipage} & \begin{minipage}[b]{\linewidth}\raggedright
Content
\end{minipage} & \begin{minipage}[b]{\linewidth}\raggedright
Forbids
\end{minipage} \\
\midrule\noalign{}
\endhead
\bottomrule\noalign{}
\endlastfoot
A1 & \textbf{Anonymity} (zeroth axiom) & Elements carry no labels; there
is no special origin, coordinate, distance, sign or type. Only relations
unchanged by every permutation of the elements are readable & Placing a
background by hand, distinguishable vertices, a ruler given from
outside, an observer standing outside \\
A2 & \textbf{Self-consistency} (first axiom) & There is no external
reference, so the system takes itself as its reference. The loop in
which the state determines itself and that state again determines itself
closes & Placing the observer or R outside the system \\
\end{longtable}
}

Provenance: A1 is Axiom 0 of the basic axiom system {[}S1{]}; A2 is A4
of the zero-closure paper {[}S3{]}.

\subsubsection{3.2 Closure and return}\label{closure-and-return}

{\def\LTcaptype{none} % do not increment counter
\begin{longtable}[]{@{}
  >{\raggedright\arraybackslash}p{(\linewidth - 6\tabcolsep) * \real{0.2500}}
  >{\raggedright\arraybackslash}p{(\linewidth - 6\tabcolsep) * \real{0.2500}}
  >{\raggedright\arraybackslash}p{(\linewidth - 6\tabcolsep) * \real{0.2500}}
  >{\raggedright\arraybackslash}p{(\linewidth - 6\tabcolsep) * \real{0.2500}}@{}}
\toprule\noalign{}
\begin{minipage}[b]{\linewidth}\raggedright
No.
\end{minipage} & \begin{minipage}[b]{\linewidth}\raggedright
Content
\end{minipage} & \begin{minipage}[b]{\linewidth}\raggedright
Explanation
\end{minipage} & \begin{minipage}[b]{\linewidth}\raggedright
Status
\end{minipage} \\
\midrule\noalign{}
\endhead
\bottomrule\noalign{}
\endlastfoot
A3 & \textbf{Closure of true squares} Σzₙ² = 0 (Z1 read at the z′ layer)
& When the whole closes consistently with itself without an external
scale, the sum of true squares (squares without conjugation) closes. The
0 is not a value but a statement that a conserved quantity can exist. By
§3.0 it is not an independent definition but the reading of the linear
closure Z1 at the layer of squared quantities & Theory. A proof as a
theorem following from self-consistency is in {[}S12{]}; that nontrivial
closure requires complex numbers is in {[}S11{]} \\
A4 & \textbf{Return} Uⁿ = I & There is no external clock, so time
evolution can only be written as a permutation among indistinguishable
elements. A permutation of finitely many indistinguishable elements has
finite order and always returns somewhere & Theory. Finite-order
resonance in {[}S7{]}, selection of roots by finite order in {[}S20{]},
the exact solution of the Z₄ phase in {[}S23{]}, the minimal exact
example in {[}S24{]} \\
\end{longtable}
}

A3 and A4 correspond to the two ways of closing. A3 is the closure of
structure, on the side where the three dimensions of space support
stationary states while remaining continuous. A4 is the closure of time,
on the side where four-dimensional spacetime becomes discrete (§3.4).
The fourteenth thought experiment {[}S19{]} treated Σx² = 0 and Uⁿ = I
not as assumptions placed at the start but as a common form appearing
when structures of different character are combined. This paper follows
that discipline and does not take either as a starting assumption: A3 is
a reading of the definition Z1 of §3.0, and A4 is a theorem.

\subsubsection{3.3 The closure equation does not fix the
scale}\label{the-closure-equation-does-not-fix-the-scale}

With zₙ = aₙ + i bₙ,

\[
\sum_n z_n^2=\sum_n(a_n^2-b_n^2)+2i\sum_n a_nb_n ,
\]

so Σzₙ² = 0 requires both Σaₙ² = Σbₙ² and Σaₙbₙ = 0. Both are
homogeneous equations and hold when every zₙ is multiplied by the same
complex number λ (Σ(λzₙ)² = λ²Σzₙ²). Scale invariance and invariance
under the global phase are built into the closure equation from the
start. The imaginary part of the closure, Σaₙbₙ = 0, is also the
condition that the generator of scale transformations vanishes {[}S1,
S6{]}.

Therefore the closure does not fix the unit of the ruler. What fixes the
ruler remains a separate question (§13).

\subsubsection{3.4 Three dimensions are continuous, four dimensions are
discrete}\label{three-dimensions-are-continuous-four-dimensions-are-discrete}

\begin{itemize}
\tightlist
\item
  \textbf{Three-dimensional space can form self-consistent states while
  remaining continuous.} Only on odd-dimensional spheres do waves
  propagate without tails and refocus sharply at the antipode (Huygens'
  principle). The spheres on which uniform translations form a group are
  only S¹ and S³. Circular waves and localized waves close as standing
  waves in space and need not advance their phase.
\item
  \textbf{Introducing time makes the phase rotate.} For self-consistency
  to keep closing, the rotation of the phase must return exactly (A4).
  With continuous time and a generic angle, the phase never returns to
  the same place and drifts quasi-periodically. To keep closing, the
  rotation angle must be a rational fraction, and time becomes discrete.
\item
  \textbf{Four-dimensional spacetime must be discrete.} For spatial
  translation and phase rotation to be compatible within one structure a
  lattice is needed, and the rotation orders compatible with a lattice
  are restricted to 1, 2, 3, 4, 6 by the crystallographic restriction.
  Four-fold and six-fold symmetries coexist on the same lattice in the
  Hurwitz lattice D₄ (24 units). There are six regular polytopes in four
  dimensions, more than the three in five or more dimensions.
\end{itemize}

\subsubsection{3.5 Other theorems and
propositions}\label{other-theorems-and-propositions}

{\def\LTcaptype{none} % do not increment counter
\begin{longtable}[]{@{}
  >{\raggedright\arraybackslash}p{(\linewidth - 6\tabcolsep) * \real{0.2500}}
  >{\raggedright\arraybackslash}p{(\linewidth - 6\tabcolsep) * \real{0.2500}}
  >{\raggedright\arraybackslash}p{(\linewidth - 6\tabcolsep) * \real{0.2500}}
  >{\raggedright\arraybackslash}p{(\linewidth - 6\tabcolsep) * \real{0.2500}}@{}}
\toprule\noalign{}
\begin{minipage}[b]{\linewidth}\raggedright
No.
\end{minipage} & \begin{minipage}[b]{\linewidth}\raggedright
Content
\end{minipage} & \begin{minipage}[b]{\linewidth}\raggedright
Status
\end{minipage} & \begin{minipage}[b]{\linewidth}\raggedright
Provenance
\end{minipage} \\
\midrule\noalign{}
\endhead
\bottomrule\noalign{}
\endlastfoot
A5 & \textbf{0 + 0 = 0 (composition)}: the closure of parts becomes the
closure of the whole. Decomposition is hard; composition is easy &
Theory & {[}S25{]} \\
A6 & \textbf{Compatibility of the norm and the true square}: a linear
map preserving both is necessarily real orthogonal & Established (§10.1)
& {[}S22{]}, Check 1 of this paper \\
\end{longtable}
}

\begin{center}\rule{0.5\linewidth}{0.5pt}\end{center}

\subsection{4. The Layer Just Below the Axioms: Big Picture, Conjugacy
and Duality, Numbers,
Projection}\label{the-layer-just-below-the-axioms-big-picture-conjugacy-and-duality-numbers-projection}

\subsubsection{4.1 Big-picture symmetries}\label{big-picture-symmetries}

These are symmetries concerning the state of the whole system rather
than individual numbers or groups. Before entering individual
symmetries, this layer is checked first.

{\def\LTcaptype{none} % do not increment counter
\begin{longtable}[]{@{}
  >{\raggedright\arraybackslash}p{(\linewidth - 8\tabcolsep) * \real{0.2000}}
  >{\raggedright\arraybackslash}p{(\linewidth - 8\tabcolsep) * \real{0.2000}}
  >{\raggedright\arraybackslash}p{(\linewidth - 8\tabcolsep) * \real{0.2000}}
  >{\raggedright\arraybackslash}p{(\linewidth - 8\tabcolsep) * \real{0.2000}}
  >{\raggedright\arraybackslash}p{(\linewidth - 8\tabcolsep) * \real{0.2000}}@{}}
\toprule\noalign{}
\begin{minipage}[b]{\linewidth}\raggedright
No.
\end{minipage} & \begin{minipage}[b]{\linewidth}\raggedright
Symmetry
\end{minipage} & \begin{minipage}[b]{\linewidth}\raggedright
Related to
\end{minipage} & \begin{minipage}[b]{\linewidth}\raggedright
Conserved quantity / fixed point
\end{minipage} & \begin{minipage}[b]{\linewidth}\raggedright
Status
\end{minipage} \\
\midrule\noalign{}
\endhead
\bottomrule\noalign{}
\endlastfoot
W1 & Time symmetry of emission and absorption, and the arrow of time &
The perfect absorber at infinity. The boundary condition of the absorber
breaks an equation that is time-symmetric in retarded and advanced waves
and selects the arrow of time (Wheeler--Feynman {[}X27{]}) & The
time-symmetric pair of fields. The record on the absorber increases
monotonically & Candidate. The absorber map is in {[}S17{]} \\
W2 & Local and global (Mach's principle, cosmological principle) & Local
R is determined by global R₀. Homogeneity and isotropy of the universe &
The representative curvature R₀ made by the eigen-oscillation of the
whole, which cancels local curvature and closes to zero & Candidate \\
W3 & The scale of observation (renormalization group), micro and macro &
α moves continuously between 128 and 137 with the energy of observation.
A closed system owns its scale. Dirac's large numbers, the resolution N₀
of the absorber & Renormalization fixed points, scale-independent ratios
& Candidate \\
W4 & Electric--magnetic duality & The Hopf fibration is the bundle of
the Dirac monopole. A single monopole quantizes charge & Electric and
magnetic charge (topological) & Candidate (set up for the first time
here) \\
W5 & Symmetry on the side of numbers (Galois symmetry) & Which rationals
are locked onto and which irrationals remain. The projection that makes
the 600-cell quasicrystal from E₈ uses the Galois conjugation sending
the golden ratio φ to −1/φ & Quantities unchanged by conjugation (field
norm and trace) & Candidate \\
W6 & Conservation of information & In a closed system information is not
lost and is recorded in the gap at the boundary & The total amount of
record (boundary area, resolution N₀) & Candidate. The minimal unit of
information is in {[}S13{]} \\
\end{longtable}
}

\subsubsection{4.2 Conjugacy and duality}\label{conjugacy-and-duality}

\textbf{Conjugacy.} Physical quantities always appear as pairs with a
fixed product: frequency and wavelength (νλ = c), position and momentum,
time and energy, mass and Compton wavelength, near and far (large and
small view angles), large and small (R and 1/R). The question at which
the programme was long stuck was what the conjugate of the amplitude A
is. As a candidate this paper names the pair of amplitude and phase. In
the exponential representation E = ln A + iφ, ln A and φ pair up as the
real and imaginary parts of the same complex number E. In quantum optics
the number N (∝ A²) and the phase are conjugate. If the phase is
periodic and becomes discrete by Uⁿ = I, its conjugate (the number) is
quantized to integers (duality of the circle U(1) and the integers ℤ).
This is a candidate for the origin of the unit of the ruler (status:
candidate). Precedents for conjugate pairs are νλ = 1 in the
dual-geometry series {[}S8{]} and the discussion of Pontryagin duality
{[}S10{]}.

\textbf{Duality.} Duality is the structural version of conjugacy. There
are two sides, such as inside and outside, spacetime and interior, front
and back; the same conserved quantity exists on both sides, but it is
called by a different name on each side. In Kaluza--Klein theory the
momentum along the fifth axis is called charge; in E = mc² the same
thing is called energy and mass; in holography the area and entropy of
the boundary are called the energy of the interior. Within this system
too, the clock of space t and the internal clock R divide one rotation
by v² + (1/γ)² = 1, and it is called velocity on one side and mass on
the other.

\textbf{Combination with anonymity.} Under anonymity, the two sides of a
duality carry no labels. Which side one reads from is only a choice of
labels, and conserved quantities called by different names on different
sides are one and the same quantity. Hence different symmetries,
conserved quantities and fixed points are often different views of one
symmetry. In the duality schema of De Haro and Butterfield {[}X12{]},
the ``common core'' shared by two models is the content of the theory,
and the interpretation maps assign meanings. In the language of this
paper, the common core corresponds to the same conserved quantity whose
name changes, and the interpretation map to the renaming. If renamings
become indistinguishable under anonymity, dual models become the same
thing. This is the claim this paper adds.

Seven groups that may merge into one are listed as candidates.

{\def\LTcaptype{none} % do not increment counter
\begin{longtable}[]{@{}
  >{\raggedright\arraybackslash}p{(\linewidth - 4\tabcolsep) * \real{0.3333}}
  >{\raggedright\arraybackslash}p{(\linewidth - 4\tabcolsep) * \real{0.3333}}
  >{\raggedright\arraybackslash}p{(\linewidth - 4\tabcolsep) * \real{0.3333}}@{}}
\toprule\noalign{}
\begin{minipage}[b]{\linewidth}\raggedright
Group
\end{minipage} & \begin{minipage}[b]{\linewidth}\raggedright
Rows that may merge
\end{minipage} & \begin{minipage}[b]{\linewidth}\raggedright
Common content
\end{minipage} \\
\midrule\noalign{}
\endhead
\bottomrule\noalign{}
\endlastfoot
K1 Allocation of one rotation & A3, S3 (Lorentz), S4 (signature), S9
(cross term 2ab), the two clocks t and R & Invariant mass = rest
frequency = rest energy \\
K2 Return of one step & A4, S2 (energy), G8 (crystallographic orders),
L3 (orders 4 and 6 of the modular group), L6 (FPUT recurrence), G14
(parity), G18 (order of the exchange) & Rotation of finite order \\
K3 Winding & G5 (winding number of the Clifford translation), S1
(momentum), G13 (Hopf winding number), W4 (magnetic charge), S6 (U(1)),
Kaluza--Klein momentum & Integer winding number \\
K4 Inside and outside & G16, W6, W1, S5, W2 & Record on the boundary =
energy of the interior \\
K5 Oneself as a fixed point & A2, L4 (renormalization fixed point), W3,
the scale freedom of A3, amplitude and phase & A point returning to
itself when the scale is changed \\
K6 Permutation of three & G11 (triality), S17 (spin and statistics), S19
(three generations), E2, S8 & Permutation of a vector and two spinors \\
K7 Labels of sides & A1, S10--S13 (C, P, T, CPT) & Which charge is
positive, which side is left, which direction is past are choices of
labels. The CPT theorem, invariance under changing all three labels at
once, is a candidate expression of the anonymity of sides \\
\end{longtable}
}

Whether a group merges is checked by whether the conserved quantities of
the group agree up to renaming alone.

\subsubsection{4.3 Dimensionless numbers attached to symmetries
(candidate axiom of magic
numbers)}\label{dimensionless-numbers-attached-to-symmetries-candidate-axiom-of-magic-numbers}

Noether's theorem says that each continuous symmetry has a conserved
quantity. The author does not regard this as an essential axiom. The
candidate axiom set up in this paper is:

\begin{quote}
An essential symmetry carries not a conserved quantity but always a
dimensionless magic number. These are the numbers hidden behind e, π and
α, and they are fixed almost entirely by number theory, group theory and
geometric structure.
\end{quote}

This system is scale invariant, and only ratios are physical (§3.3).
What a symmetry fixes is not a dimensional conserved quantity but a
dimensionless number. A conserved quantity is only that number read out
after multiplying by a ruler. In a world of ratios only, every ratio is
preserved however the ruler is changed, and there are no special values.
What does not vanish under a change of scale is only the pure numbers
determined by structure. Integer numbers (orders, winding numbers,
counts) select allowed values discretely from continuous ratios and
bring in discreteness. Numbers that determine groups, lattices and
geometries (π, φ, E₈'s 240 and 248) determine the form of relations and
bring in structure. Scale invariance cannot break itself, so in a system
that brings in no ruler from outside, the only source of discreteness
and structure left is magic numbers.

Known examples abound. The nuclear magic numbers 2, 8, 20, 28, 50, 82,
126 come from adding a spin--orbit term to the SU(3) symmetry of the
three-dimensional harmonic oscillator (2, 8, 20, 40, 70, 112) {[}X30{]}.
Universal constants at renormalization fixed points (the Feigenbaum
constants, the dimension of the devil's staircase) are determined by
structure alone, independently of the details of the system {[}X23{]}.
E₈ has many numbers, in proportion to the many symmetries imposing
constraints (§9).

When different symmetries require nearby numbers, the observed value
does not settle on any one of them; it is pulled in and oscillates
metastably between them. Observed constants are close to, but not equal
to, structural numbers. Eddington demanded agreement by claiming that
α⁻¹ is the integer 136, later 137, while the mainstream discarded the
structural numbers themselves because of the gap {[}X18{]}. This paper
takes neither side and treats the gap not as an error but as content to
be read. This relation is why a world that is uncertain and discrete
looks continuous and deterministic. Because structural numbers give the
points of locking, jumps are observed discretely. Because several
numbers compete, observed values drift continuously with scale (the
running of α). Quasi-periodic motion that is not locked takes
exponentially long to break down (KAM and Nekhoroshev), and meanwhile it
looks like a fixed orbit.

\subsubsection{4.4 The projection
symmetry}\label{the-projection-symmetry}

{\def\LTcaptype{none} % do not increment counter
\begin{longtable}[]{@{}
  >{\raggedright\arraybackslash}p{(\linewidth - 8\tabcolsep) * \real{0.2000}}
  >{\raggedright\arraybackslash}p{(\linewidth - 8\tabcolsep) * \real{0.2000}}
  >{\raggedright\arraybackslash}p{(\linewidth - 8\tabcolsep) * \real{0.2000}}
  >{\raggedright\arraybackslash}p{(\linewidth - 8\tabcolsep) * \real{0.2000}}
  >{\raggedright\arraybackslash}p{(\linewidth - 8\tabcolsep) * \real{0.2000}}@{}}
\toprule\noalign{}
\begin{minipage}[b]{\linewidth}\raggedright
No.
\end{minipage} & \begin{minipage}[b]{\linewidth}\raggedright
Symmetry
\end{minipage} & \begin{minipage}[b]{\linewidth}\raggedright
Related to
\end{minipage} & \begin{minipage}[b]{\linewidth}\raggedright
Conserved quantity / fixed point
\end{minipage} & \begin{minipage}[b]{\linewidth}\raggedright
Status
\end{minipage} \\
\midrule\noalign{}
\endhead
\bottomrule\noalign{}
\endlastfoot
P1 & Full, partial and zero projection of an operation interchange with
the projection method & The map between state space and spacetime,
uncertainty, discreteness, the Born rule, the exchange ratio & The ratio
of projection (a fixed point of locking if stable) & Hypothesis. Axiom
0.7 of the basic axiom system {[}S1{]}, the anonymous partial projection
paper {[}S2{]} \\
\end{longtable}
}

The same state and the same operation can look like a full projection, a
partial projection or a zero projection depending on the projection
through which they are viewed. A complex number looks like a full
projection in one projection (it maps onto all values) and like a zero
projection in another (multiplying by i leaves zero projection onto the
original direction). The two are not contradictory; they are different
projections.

{\def\LTcaptype{none} % do not increment counter
\begin{longtable}[]{@{}
  >{\raggedright\arraybackslash}p{(\linewidth - 4\tabcolsep) * \real{0.3333}}
  >{\raggedright\arraybackslash}p{(\linewidth - 4\tabcolsep) * \real{0.3333}}
  >{\raggedright\arraybackslash}p{(\linewidth - 4\tabcolsep) * \real{0.3333}}@{}}
\toprule\noalign{}
\begin{minipage}[b]{\linewidth}\raggedright
How it looks
\end{minipage} & \begin{minipage}[b]{\linewidth}\raggedright
Examples
\end{minipage} & \begin{minipage}[b]{\linewidth}\raggedright
What comes out
\end{minipage} \\
\midrule\noalign{}
\endhead
\bottomrule\noalign{}
\endlastfoot
As a full projection & The norm (always positive, never 0),
magnification and reduction & Continuity, definiteness, determinism \\
As a partial projection & Rotation by an angle θ (cos θ remains in the
original direction), measurement of a superposition (probability cos²θ),
the right exchange R = cos²θ & Probability, locking ratios \\
As a zero projection & Multiplication by i, the antisymmetric part of a
product (commutator), measurement of an orthogonal state & Uncertainty
(bounded below by the commutator, {[}x, p{]} = iħ), discreteness (i⁴ =
1) \\
\end{longtable}
}

The map from the state space \{zₙ\} to spacetime corresponds to viewing
the state space through some projection. Until the projection method is
chosen, what is readable and what is not is undetermined. \textbf{This
is the root difficulty of the map between state space and spacetime.}
Since spacetime is relatively stably readable, the projection must be at
least a partial projection. With a zero projection spacetime would not
appear; with a full projection neither the unreadable interior (RQ₁Q₂Q₃)
nor uncertainty nor discreteness would remain. With a partial
projection, the projected part becomes readable spacetime \{xyzt\}, the
part escaping orthogonally becomes the unreadable interior and phase,
and uncertainty and discreteness come from the escaped part. For
spacetime to be stably readable the ratio of projection must also be
stable, and a stable ratio should be locked onto some structural number.
The right exchange R = cos²θ ≈ 0.7 has the form of a partial-projection
ratio, but whether it is the candidate is unknown (candidate).

\begin{center}\rule{0.5\linewidth}{0.5pt}\end{center}

\subsection{5. Symmetries Required by the Standard Model, General
Relativity and
E₈}\label{symmetries-required-by-the-standard-model-general-relativity-and-eux2088}

External standards are listed with their candidate correspondences in
this system.

{\def\LTcaptype{none} % do not increment counter
\begin{longtable}[]{@{}
  >{\raggedright\arraybackslash}p{(\linewidth - 10\tabcolsep) * \real{0.1667}}
  >{\raggedright\arraybackslash}p{(\linewidth - 10\tabcolsep) * \real{0.1667}}
  >{\raggedright\arraybackslash}p{(\linewidth - 10\tabcolsep) * \real{0.1667}}
  >{\raggedright\arraybackslash}p{(\linewidth - 10\tabcolsep) * \real{0.1667}}
  >{\raggedright\arraybackslash}p{(\linewidth - 10\tabcolsep) * \real{0.1667}}
  >{\raggedright\arraybackslash}p{(\linewidth - 10\tabcolsep) * \real{0.1667}}@{}}
\toprule\noalign{}
\begin{minipage}[b]{\linewidth}\raggedright
No.
\end{minipage} & \begin{minipage}[b]{\linewidth}\raggedright
Symmetry
\end{minipage} & \begin{minipage}[b]{\linewidth}\raggedright
Conserved quantity
\end{minipage} & \begin{minipage}[b]{\linewidth}\raggedright
Noether's theorem
\end{minipage} & \begin{minipage}[b]{\linewidth}\raggedright
Correspondence in this system
\end{minipage} & \begin{minipage}[b]{\linewidth}\raggedright
Status
\end{minipage} \\
\midrule\noalign{}
\endhead
\bottomrule\noalign{}
\endlastfoot
S1 & Spatial translation & Momentum & Holds; not localizable in gravity
& Left Clifford translation (fixed-point free; winding number and
one-step angle are gauge invariants) & Candidate \\
S2 & Time translation & Energy & Holds classically & Eigenphase of Uⁿ =
I & Candidate \\
S3 & Lorentz transformation & Angular momentum, boost charges & Holds &
Decomposition z² = (a² − b²) + i·2ab, v² + (1/γ)² = 1, additivity of
rapidity & Identity is theory {[}S16, S21{]} \\
S4 & Signature of spacetime & ― & ― & Relocating terms by multiplying a
coordinate by i; groups of different signature all lie in the same O(n,
ℂ) & Established (§10.1) \\
S5 & General coordinate transformation, equivalence principle & No local
conserved quantity (identity by the second theorem) & Doubtful & Stretch
γ by a localized wave, deficit angles of view-angle triangles &
Candidate {[}S15{]} \\
S6 & U(1)\_Y and U(1)\_EM & Charge & Holds & Hopf winding number, the
charge formula of sign vectors Q\_em = −(1/6)(Q₁ + Q₂ + Q₃) − t/2 &
Candidate \\
S7 & SU(2)\_L & Weak isospin & Spontaneously broken & One of SU(2)\_L ×
SU(2)\_R, simultaneous flip of (t, R) & Candidate \\
S8 & SU(3) & Colour charge & Holds & G₂ ⊃ SU(3) fixed by triality, or
the 3 + 2 decomposition {[}S3, S18{]} & Candidate \\
S9 & Electroweak symmetry breaking & ― & Broken & Pinned cross term 2ab
(spin 0) & Candidate \\
S10--S13 & C, P, T, CPT & ― & Discrete & Sign flip of Q, quaternion
conjugation q ↦ q̄ (exchange of left and right multiplication), complex
conjugation & Candidate \\
S14 & CP violation & ― & Broken & Mismatch between the exchange rotation
and the lattice & Candidate \\
S15 & Baryon number, lepton number & B, L & B + L anomalous at the
quantum level & Components of sign vectors & Open \\
S16 & Chiral symmetry & Axial charge & Broken by quantum anomaly & Left
and right multiplication {[}S4{]} & Candidate \\
S17 & Spin and statistics & Type of statistics & Follows from Lorentz
invariance & Parity (−1)ⁿ & Candidate \\
S18 & Gauge anomaly cancellation & Sum of charges per generation is 0 &
Consistency condition & Charges of sign vectors come out without excess
or deficit {[}S4{]} & Candidate \\
S19 & Three generations & ― & Reason unknown & Triality (order 3) &
Candidate \\
S20 & Scale and conformal symmetry & Dilation charge & Broken by mass
and quantum anomaly & Σz² = 0 is trivially scale invariant & Theory
(§3.3) \\
\end{longtable}
}

{\def\LTcaptype{none} % do not increment counter
\begin{longtable}[]{@{}
  >{\raggedright\arraybackslash}p{(\linewidth - 6\tabcolsep) * \real{0.2500}}
  >{\raggedright\arraybackslash}p{(\linewidth - 6\tabcolsep) * \real{0.2500}}
  >{\raggedright\arraybackslash}p{(\linewidth - 6\tabcolsep) * \real{0.2500}}
  >{\raggedright\arraybackslash}p{(\linewidth - 6\tabcolsep) * \real{0.2500}}@{}}
\toprule\noalign{}
\begin{minipage}[b]{\linewidth}\raggedright
No.
\end{minipage} & \begin{minipage}[b]{\linewidth}\raggedright
Where E₈ appears
\end{minipage} & \begin{minipage}[b]{\linewidth}\raggedright
Correspondence in this system
\end{minipage} & \begin{minipage}[b]{\linewidth}\raggedright
Status
\end{minipage} \\
\midrule\noalign{}
\endhead
\bottomrule\noalign{}
\endlastfoot
E1 & Gluing of the readable side and the interior & Pairing two D₄
lattices in the same coset gives the 240 roots of E₈ & Established
(§10.2); precedent {[}S2{]} \\
E2 & E₈ ⊃ SO(16) ⊃ D₄ × D₄ & Triality of the pair preserves the roots of
E₈ & Established (§10.2) \\
E3 & Grand-unification chain E₈ ⊃ E₆ ⊃ SO(10) ⊃ SU(5) & ― & Open \\
E4 & The spinor 16 of SO(10) & Same number as the 16 even sign vectors
tRQ₁Q₂Q₃ & Candidate (set up for the first time here) \\
E5 & Heterotic string, Lisi's E₈ theory & ― & Reference \\
E6 & Contact between E₈ and 137 & ― & Open \\
\end{longtable}
}

\begin{center}\rule{0.5\linewidth}{0.5pt}\end{center}

\subsection{6. Symmetries of the Algebra and Geometry of This
System}\label{symmetries-of-the-algebra-and-geometry-of-this-system}

{\def\LTcaptype{none} % do not increment counter
\begin{longtable}[]{@{}
  >{\raggedright\arraybackslash}p{(\linewidth - 8\tabcolsep) * \real{0.2000}}
  >{\raggedright\arraybackslash}p{(\linewidth - 8\tabcolsep) * \real{0.2000}}
  >{\raggedright\arraybackslash}p{(\linewidth - 8\tabcolsep) * \real{0.2000}}
  >{\raggedright\arraybackslash}p{(\linewidth - 8\tabcolsep) * \real{0.2000}}
  >{\raggedright\arraybackslash}p{(\linewidth - 8\tabcolsep) * \real{0.2000}}@{}}
\toprule\noalign{}
\begin{minipage}[b]{\linewidth}\raggedright
No.
\end{minipage} & \begin{minipage}[b]{\linewidth}\raggedright
Symmetry
\end{minipage} & \begin{minipage}[b]{\linewidth}\raggedright
Related to
\end{minipage} & \begin{minipage}[b]{\linewidth}\raggedright
Status
\end{minipage} & \begin{minipage}[b]{\linewidth}\raggedright
Provenance
\end{minipage} \\
\midrule\noalign{}
\endhead
\bottomrule\noalign{}
\endlastfoot
G1 & Composition algebras (dimensions 1, 2, 4, 8) & The eight directions
& Established & Hurwitz's theorem \\
G2 & S³ = SU(2), uniform translations form a group & The three
dimensions of space (only S¹ and S³; S⁷ violates associativity) &
Established & Adams' theorem {[}S21{]} \\
G3 & Huygens' principle & Stationarity of localized waves & Established
& {[}S21{]} \\
G4 & Mass gap from curvature & The −1 in the S³ eigenvalues (n + 1)² − 1
& Established & {[}S21{]} \\
G5 & Left multiplication (Clifford translation) & An equally spaced
scale unaffected by the distortion of space (gauge) & Theory &
{[}S21{]} \\
G6 & Right multiplication (array of component rotations, exchange) &
Interaction, R = cos²θ ≈ 0.7 & Theory & {[}S7, S27{]} \\
G7 & Continuous rotations in three dimensions & Angular momentum; the
centre θ₀ is a fixed point & Established & ― \\
G8 & Crystallographic symmetry in four dimensions & Discreteness of
spacetime; orders 1, 2, 3, 4, 6 (see the open item in §11.3) & Theory &
§3.4 \\
G9 & The 24 units of the R = 3 system & Multiplication from left or
right preserves the lattice H and D₄; the ball centre (½,½,½,½) is one
of the units & Established & Checks of this paper \\
G10 & Odd R & Vertices lie on a cell-centred hypersphere only for odd R;
R = 3 gives 137 & Established (count) & {[}S9, S31, S32{]} \\
G11 & Triality & One vector set and two spinor sets, candidate for three
generations & Established (math), candidate (physics) & §10.2,
{[}X28{]} \\
G12 & Quaternion form of triality & q ↦ a q a & Established & §10.2 (set
up for the first time here) \\
G13 & Hopf fibration & Charge (winding number) & Theory & {[}S21{]} \\
G14 & Parity (−1)ⁿ, antipode & Fermions and bosons & Theory &
{[}S21{]} \\
G15 & Sign vectors tRQ₁Q₂Q₃ & 16 even ones form a generation; xyz
carries spin & Theory & This paper (study of 2026-10-16) \\
G16 & Duality of inside and outside & Cross-polytope and hypercube,
particle and absorber & Candidate & The mathematics of the
cross-polytope--hypercube duality is in {[}S28{]} \\
G17 & Central projection & R is unreadable; the origin is shared &
Theory & Axiom 3 of the basic axiom system {[}S1{]}, chapter 11 of the
fourteenth {[}S19{]} \\
G18 & Exchange rotation 23π/124 & Finite order, not crystallographic;
not a symmetry of the lattice & Established (not a lattice symmetry) &
{[}S7, S19{]} \\
\end{longtable}
}

\begin{center}\rule{0.5\linewidth}{0.5pt}\end{center}

\subsection{7. Symmetries of Locking and Metastable
Oscillation}\label{symmetries-of-locking-and-metastable-oscillation}

When nearby discrete values are required by different symmetries, the
system may not lock onto any of them and may instead oscillate
metastably or asymptotically with an extremely long period. Symmetries
found by complex systems and nonlinear dynamics are listed with their
correspondences in this system.

{\def\LTcaptype{none} % do not increment counter
\begin{longtable}[]{@{}
  >{\raggedright\arraybackslash}p{(\linewidth - 8\tabcolsep) * \real{0.2000}}
  >{\raggedright\arraybackslash}p{(\linewidth - 8\tabcolsep) * \real{0.2000}}
  >{\raggedright\arraybackslash}p{(\linewidth - 8\tabcolsep) * \real{0.2000}}
  >{\raggedright\arraybackslash}p{(\linewidth - 8\tabcolsep) * \real{0.2000}}
  >{\raggedright\arraybackslash}p{(\linewidth - 8\tabcolsep) * \real{0.2000}}@{}}
\toprule\noalign{}
\begin{minipage}[b]{\linewidth}\raggedright
No.
\end{minipage} & \begin{minipage}[b]{\linewidth}\raggedright
Known structure
\end{minipage} & \begin{minipage}[b]{\linewidth}\raggedright
Content
\end{minipage} & \begin{minipage}[b]{\linewidth}\raggedright
Correspondence in this system
\end{minipage} & \begin{minipage}[b]{\linewidth}\raggedright
Status
\end{minipage} \\
\midrule\noalign{}
\endhead
\bottomrule\noalign{}
\endlastfoot
L1 & Arnold tongues, circle map & Frequency ratios lock onto nearby
rationals; the width narrows as the denominator grows & Locking between
the exchange rotation and the lattice orders; measured in the model in
{[}S26{]} & Established (math), candidate (correspondence) \\
L2 & Devil's staircase & At criticality a step stands at every rational;
the remainder is a Cantor set of dimension about 0.868 {[}X23{]} & The
structure, as a set, of being near but not equal & Established,
candidate \\
L3 & Farey tree, modular group SL(2, ℤ) & Steps are ordered by mediants
(p + p′)/(q + q′) {[}X24{]}. The modular group is generated by elements
of order 4 and 6 & The crystallographic orders Z₄ and Z₆ enter the
symmetry of the locking hierarchy. In the quantum Hall effect the
subgroup Γ₀(2) gives the selection rule of transitions {[}X25{]} &
Established, candidate \\
L4 & Golden ratio, renormalization fixed point & The ratio least prone
to locking; scale symmetry of self-similarity & Limit of metastability &
Established, candidate \\
L5 & KAM theorem, Nekhoroshev stability & Quasi-periodicity survives;
even when it breaks, the time to break down is exponentially long &
Asymptotic oscillation with an extremely long period & Established,
candidate \\
L6 & FPUT recurrence & No approach to thermal equilibrium; near return
after a long time & Approximate return of Uⁿ = I & Established,
candidate \\
L7 & Kuramoto model & Transition between synchrony and asynchrony & How
the phases of many zₙ align; the single-wave equation is of Kuramoto
type {[}S22{]} & Established, candidate \\
L8 & Quasicrystals & Projecting E₈ to four dimensions gives a
quasicrystal with the 600-cell symmetry H₄, with the golden ratio
entering {[}X26{]} & A symmetry incompatible with the lattice appears as
an order that is near but not equal & Established, candidate \\
L9 & Theta function of E₈ & A modular form (Eisenstein series of weight
4) & Contact between E₈ and the modular group & Established,
candidate \\
\end{longtable}
}

Locking appears not as a conserved quantity of a flow but as a fixed
point (the lock) and quasi-oscillation around it. The symmetries of the
devil's staircase, the Farey tree and renormalization fixed points are
discrete hierarchies and scale self-similarity, outside Noether's
theorem.

\begin{center}\rule{0.5\linewidth}{0.5pt}\end{center}

\subsection{8. Where Noether's Theorem Fails or Is
Doubtful}\label{where-noethers-theorem-fails-or-is-doubtful}

\begin{enumerate}
\def\labelenumi{\arabic{enumi}.}
\tightlist
\item
  \textbf{Discrete symmetries} (A1, A4, G8, G11, G14, G15): there is no
  current. What is conserved is an eigenvalue (a phase or a sign).
\item
  \textbf{Fixed points} (A2): self-consistency is a fixed point rather
  than a symmetry; the conserved quantity and the fixed point are the
  same thing.
\item
  \textbf{Topological conservation} (G13, W4): winding numbers are
  conserved by topological structure rather than by symmetry.
\item
  \textbf{Scale} (A3, S20): there is a symmetry, yet the closure does
  not fix the ruler.
\item
  \textbf{Gravity} (S5): local stretching breaks translation symmetry
  and energy cannot be localized. By Noether's second theorem
  conservation laws become identities holding identically (``strong
  laws'') {[}X8{]}.
\item
  \textbf{Quantum anomalies} (S15, S16, S20): the axial U(1), B + L and
  scale symmetry are conserved classically but broken by quantization
  {[}X9{]}.
\item
  \textbf{Big-picture symmetries} (W1, W4, W5): the arrow of time is
  broken by a boundary condition, and electric--magnetic duality and
  Galois symmetry lie outside Noether's theorem.
\item
  \textbf{Mismatch of the exchange} (G18): there is no exact conserved
  quantity, only locking and quasi-oscillation.
\item
  \textbf{Locking symmetries} (§7): they appear as fixed points and
  quasi-oscillation around them.
\end{enumerate}

All of these lie on the side of the view of §4.3 that what is attached
to a symmetry is not a conserved quantity but a number.

\begin{center}\rule{0.5\linewidth}{0.5pt}\end{center}

\subsection{9. Catalogue of Magic
Numbers}\label{catalogue-of-magic-numbers}

Numbers that look like magic numbers are divided into four systems by
their origin. When a new number appears, first decide its system.
Mistaking a derived number for an integer or irrational one leads into
the trap of demanding agreement.

\subsubsection{9.1 Integer system (numbers required by
structure)}\label{integer-system-numbers-required-by-structure}

\textbf{Easily overlooked integers}: 0 (closure, a statement of
existence), 1 (one closure, v² + (1/γ)² = 1), 2 (one complex number
cannot close; the minimal system is a two-body system), 3 (the dimension
of space, triality, three generations), 4 (the dimension of spacetime,
quaternions), 8 (the upper limit of composition algebras, the eight
directions), n (the order of Uⁿ = I).

\textbf{Integers required as sequences} (mostly from number theory,
group theory and geometry):

{\def\LTcaptype{none} % do not increment counter
\begin{longtable}[]{@{}
  >{\raggedright\arraybackslash}p{(\linewidth - 4\tabcolsep) * \real{0.3333}}
  >{\raggedright\arraybackslash}p{(\linewidth - 4\tabcolsep) * \real{0.3333}}
  >{\raggedright\arraybackslash}p{(\linewidth - 4\tabcolsep) * \real{0.3333}}@{}}
\toprule\noalign{}
\begin{minipage}[b]{\linewidth}\raggedright
Sequence
\end{minipage} & \begin{minipage}[b]{\linewidth}\raggedright
Origin
\end{minipage} & \begin{minipage}[b]{\linewidth}\raggedright
Status
\end{minipage} \\
\midrule\noalign{}
\endhead
\bottomrule\noalign{}
\endlastfoot
1, 2, 4, 8 & Dimensions of composition algebras & Established \\
1, 3, 7 & Dimensions of parallelizable spheres (S¹ and S³ form groups) &
Established \\
1, 3, 7, 15 & Dimensions of the total spaces of Hopf fibrations &
Established \\
1, 2, 3, 4, 6 & Rotation orders compatible with a lattice &
Established \\
5, 6, 3 & Numbers of regular polyhedra and polytopes (three, four, five
or more dimensions) & Established \\
2, 6, 12, 24, 240, 196560 & Kissing numbers (dimensions 1, 2, 3, 4, 8,
24) & Established \\
24, 48, 120 & Orders of the binary tetrahedral, octahedral and
icosahedral groups & Established \\
14, 52, 78, 133, 248 & Dimensions of the exceptional Lie groups G₂, F₄,
E₆, E₇, E₈ & Established \\
2, 8, 20, 28, 50, 82, 126 & Nuclear magic numbers & Established \\
1, 137, 1545, 7281 & Numbers of hypercubes inscribed in a cell-centred
hypersphere (R = 1, 3, 5, 7) & Established (count) {[}S9{]} \\
1, 2, 3, 7, 11, 19, 43, 67, 163 & Heegner numbers & Established \\
8 & Bott periodicity & Established \\
\end{longtable}
}

\textbf{Integers of E₈ and related lattices} ({[}X29{]}): rank 8,
dimension 248, 240 roots, Coxeter number 30 (also 30 for the 600-cell
group H₄), Weyl group order 696,729,600, determinant 1. The exponents of
E₈, 1, 7, 11, 13, 17, 19, 23, 29, are the integers below 30 coprime to
30; their number is φ(30) = 8 = rank, and their sum 120 is the number of
positive roots, equal to the number of vertices of the 600-cell. The
binary icosahedral group (the 120 vertices of the 600-cell) corresponds
to E₈ by the McKay correspondence. For D₄: 24 roots, Weyl group order
192, 1152 automorphisms (§10.2).

\textbf{Ratios of integers}: 1/2 of spin and of Dirac's condition,
charges 1/3 and 2/3, 16 per generation, 12 gauge particles (8 + 3 + 1).

\subsubsection{9.2 Coupling-convention system (numbers set by unit
conventions)}\label{coupling-convention-system-numbers-set-by-unit-conventions}

G = c = q₀ = 1, q₀ = e/3, J = ħ/2, natural units ħ = c = 1. These are
not magic numbers but numbers set to 1 or 3 by the choice of units. Only
dimensional constants can be set to 1 by a unit convention;
dimensionless numbers (α and others) cannot be changed by the choice of
units.

\subsubsection{9.3 Irrational system}\label{irrational-system}

\begin{itemize}
\tightlist
\item
  \textbf{Geometry (space)}: π, √2, √3, the golden ratio φ, the dihedral
  angle of the regular tetrahedron arccos(1/3) ≈ 70.53° (five of them
  fall short of 360° by 7.356°), the complement of the magic angle
  35.264°, the solid angle 4π.
\item
  \textbf{Evolution (phase, time)}: e (the exponential map of a group),
  2π/n (one step of Uⁿ = I), e\^{}\{iπ\} = −1 (connecting geometry and
  evolution), rapidity.
\item
  \textbf{Statistics (an easily overlooked system)}: √(2π), Euler's
  constant, ln 2, ζ(3), π⁴/15, the constants of Wien's displacement law,
  Onsager's critical temperature 2/ln(1 + √2), percolation thresholds,
  the Feigenbaum constants, the dimension 0.868 of the devil's
  staircase, the golden-mean critical constant of the circle map,
  Khinchin's constant, Lévy's constant π²/(12 ln 2), the critical
  coupling of the Kuramoto model, critical exponents.
\end{itemize}

\subsubsection{9.4 Derived system (numbers that look like magic numbers
but are derived from other
numbers)}\label{derived-system-numbers-that-look-like-magic-numbers-but-are-derived-from-other-numbers}

α⁻¹ = 137.035999\ldots{} (the electromagnetic coupling, running with
energy), the Thomson cross-section (8π/3) r\_e², the first term of the
anomalous magnetic moment α/2π, fine-structure splitting ∝ α², the
Rydberg constant α² m\_e c/(2h), the Bohr radius λ̄\_C/α, sin²θ\_W, the
proton--electron mass ratio, the strong coupling, the gravitational
coupling. These should be assembled by the readout map from numbers of
the integer and irrational systems, and they appear as cross-sections
and interaction coefficients.

\subsubsection{9.5 Pairs that are near but not
equal}\label{pairs-that-are-near-but-not-equal}

{\def\LTcaptype{none} % do not increment counter
\begin{longtable}[]{@{}
  >{\raggedright\arraybackslash}p{(\linewidth - 6\tabcolsep) * \real{0.2500}}
  >{\raggedright\arraybackslash}p{(\linewidth - 6\tabcolsep) * \real{0.2500}}
  >{\raggedright\arraybackslash}p{(\linewidth - 6\tabcolsep) * \real{0.2500}}
  >{\raggedright\arraybackslash}p{(\linewidth - 6\tabcolsep) * \real{0.2500}}@{}}
\toprule\noalign{}
\begin{minipage}[b]{\linewidth}\raggedright
Structural number
\end{minipage} & \begin{minipage}[b]{\linewidth}\raggedright
Observed or nearby number
\end{minipage} & \begin{minipage}[b]{\linewidth}\raggedright
Explanation
\end{minipage} & \begin{minipage}[b]{\linewidth}\raggedright
Status
\end{minipage} \\
\midrule\noalign{}
\endhead
\bottomrule\noalign{}
\endlastfoot
137 & α⁻¹ = 137.036 & Unexplained. Candidate for locking and
quasi-oscillation among several structural numbers & Candidate \\
640320³ + 744 & e\^{}\{π√163\} (smaller by 7.5 × 10⁻¹³) & Explained by
class number 1 and the expansion of the j-function {[}X32{]} &
Established \\
196883 + 1 & The coefficient 196884 of the j-function & The nearness
revealed a hidden symmetry (the Monster) and became a theorem {[}X31{]}
& Established \\
2/3 & Koide's value & Unexplained & Candidate \\
Lattice orders 4 and 6 & Order of the exchange & Not crystallographic;
locking in the fourteenth {[}S19{]} & Theory \\
\end{longtable}
}

e\^{}\{π√163\} and moonshine are examples where ``near but not equal''
comes from the structure of number theory and group theory itself, and
they are precedents for a method of finding hidden symmetries from near
coincidences.

\begin{center}\rule{0.5\linewidth}{0.5pt}\end{center}

\subsection{10. Confirmed Results}\label{confirmed-results}

The checks in this section were made only with the mathematics of
objects already in the theory, without using the readout map (which is
open). No parts were added. Programs and results are in the folder of
Appendix A.

\subsubsection{10.1 Check 1: true square and
norm}\label{check-1-true-square-and-norm}

\textbf{Proposition.} A linear map T that preserves both the sum of true
squares zᵀz and the norm z†z is necessarily a real orthogonal map.

\textbf{Proof.} From TᵀT = I and T†T = I we get Tᵀ = T†, that is, T̄ = T,
so T is real. Being real with TᵀT = I, it is orthogonal. For generators,
the intersection of complex antisymmetric and anti-Hermitian matrices is
the real antisymmetric matrices. □

Numerically, in eight dimensions, complex orthogonal maps preserved the
true square (change 3.6 × 10⁻¹⁴) and changed the norm (change 100),
unitary maps preserved the norm (1.8 × 10⁻¹⁵) and changed the true
square (2.9), and real orthogonal maps preserved both (10⁻¹⁵). The
dimension of the intersection of generators was 28, equal to the
dimension of real antisymmetric matrices in eight dimensions.

\textbf{Relocating terms.} With D = diag(1, \ldots, i, \ldots), which
multiplies coordinates by i, an element Λ of the group O(p, q) of
signature (4,0), (3,1) or (2,2) is mapped to DΛD⁻¹, and all of these
became elements of the same O(4, ℂ) (error 10⁻¹⁶). Exchanging space and
time corresponds to a different real slice of the same complex quadric.

\subsubsection{10.2 Check 2: triality of
D₄}\label{check-2-triality-of-dux2084}

{\def\LTcaptype{none} % do not increment counter
\begin{longtable}[]{@{}
  >{\raggedright\arraybackslash}p{(\linewidth - 2\tabcolsep) * \real{0.5000}}
  >{\raggedright\arraybackslash}p{(\linewidth - 2\tabcolsep) * \real{0.5000}}@{}}
\toprule\noalign{}
\begin{minipage}[b]{\linewidth}\raggedright
What was checked
\end{minipage} & \begin{minipage}[b]{\linewidth}\raggedright
Result
\end{minipage} \\
\midrule\noalign{}
\endhead
\bottomrule\noalign{}
\endlastfoot
Automorphisms of the D₄ root system & 1152 (the order of W(F₄)) \\
Action on the three cosets of D₄*/D₄: v (odd class), s and c (two
half-integer classes) & Each of the 6 permutations occurs 192 times;
triality is present in full as S₃ \\
An element T of order 3 cycling v → s → c & Exists \\
The 24 Hurwitz units & Split into 8 each in v, s, c; T cycles the three
octets and preserves the set of 24 \\
The pair (T, T) & Preserves the 240 roots of E₈ \\
The centre (½,½,½,½) of the ball in which 137 was counted & In class s;
T sends it to class c, T² to class v \\
Quaternion form of T & \textbf{q ↦ a q a}, a = (−1 + i + j + k)/2 (a
unit of order 3); the operation of multiplying by the same unit from
left and right \\
\end{longtable}
}

The 24 units of the equally spaced gauge of the R = 3 system split into
one vector set and two spinor sets, and triality exchanges them by left
and right multiplication. The subgroup fixed by triality is G₂, and G₂
contains SU(3) (a mathematical fact). Whether this corresponds to
physical colour or to three generations is a candidate.

\begin{center}\rule{0.5\linewidth}{0.5pt}\end{center}

\subsection{11. Correspondence with Published
Papers}\label{correspondence-with-published-papers}

The 156 published papers (13 series) were mapped to the rows of the
catalogue. 81 rows have a published standard; the 15 rows to which no
paper corresponds are W4, K3, K6, K7, S13, S15, E4, E5, G12, C3, C4, L4,
L5, L8 and L9. These 15 rows are supported by external literature or by
the checks of this paper. The full mapping table is in Appendix B.

\subsubsection{11.1 Resetting the axiom
numbering}\label{resetting-the-axiom-numbering}

{\def\LTcaptype{none} % do not increment counter
\begin{longtable}[]{@{}
  >{\raggedright\arraybackslash}p{(\linewidth - 4\tabcolsep) * \real{0.3333}}
  >{\raggedright\arraybackslash}p{(\linewidth - 4\tabcolsep) * \real{0.3333}}
  >{\raggedright\arraybackslash}p{(\linewidth - 4\tabcolsep) * \real{0.3333}}@{}}
\toprule\noalign{}
\begin{minipage}[b]{\linewidth}\raggedright
Document
\end{minipage} & \begin{minipage}[b]{\linewidth}\raggedright
Old number
\end{minipage} & \begin{minipage}[b]{\linewidth}\raggedright
Number in this paper
\end{minipage} \\
\midrule\noalign{}
\endhead
\bottomrule\noalign{}
\endlastfoot
Basic axiom system v9.1 {[}S1{]} & Axiom 0 anonymity & A1 \\
Same & Axiom 0.5 scale anonymity & Consequence of A3, S20 \\
Same & Axiom 0.6 dimension-wise central-projection readout & No
corresponding row (a readout rule for curvature correction) \\
Same & Axiom 0.7 existence of anonymous partial projection & P1 \\
Same & Axiom 1 all-positive-sign zero closure & A3 (moved from axiom to
theorem) \\
Same & Axiom 2 nontrivial existence & Premise of A3 \\
Same & Axiom 3 degeneracy of the projection axis and imaginary-direction
readout & G17 \\
Same & Axiom 15 discreteness of observed phase & A4, G13 \\
Zero closure, finite order, self-consistent geometry {[}S3{]} & A1
complex zero closure & A3 \\
Same & A2 finite order & A4 \\
Same & A3 simplex consistency & No corresponding row (later found to
impose no constraint {[}S12{]}) \\
Same & A4 self-consistency & A2 \\
Earlier N-body series & Axiom 1, first axiom (zero closure) & A3 \\
Research project design {[}S14{]} & First-principle candidate (anonymous
distinguishable states) & A1 \\
First and second thought experiments {[}S15, S16{]} & P1 anonymity, A1 a
conserved readout & A1, A3 and K1 \\
Dual-geometry series & νλ = 1, zero-point ½, one asymmetric bit & The
conjugate pairs of §4.2, the integer ratio ½, S12 \\
\end{longtable}
}

\subsubsection{11.2 Revision of a claim}\label{revision-of-a-claim}

Papers 7 and 9 and the supplement of the BH thermodynamics programme
{[}S29, S30{]} took the locking side, filling the gap between 137 and
α⁻¹ with a closed identity. This paper revises that to the position that
reads the gap itself as content (§4.3, §9.5). The earlier papers are
cited as the record of the position before revision.

\subsubsection{11.3 Interpretations of numbers kept as open
problems}\label{interpretations-of-numbers-kept-as-open-problems}

For the following disagreements, no single interpretation is taken as
correct. That several interpretations are possible is the present
answer, and these remain as problems.

{\def\LTcaptype{none} % do not increment counter
\begin{longtable}[]{@{}
  >{\raggedright\arraybackslash}p{(\linewidth - 2\tabcolsep) * \real{0.5000}}
  >{\raggedright\arraybackslash}p{(\linewidth - 2\tabcolsep) * \real{0.5000}}@{}}
\toprule\noalign{}
\begin{minipage}[b]{\linewidth}\raggedright
Item
\end{minipage} & \begin{minipage}[b]{\linewidth}\raggedright
Interpretations
\end{minipage} \\
\midrule\noalign{}
\endhead
\bottomrule\noalign{}
\endlastfoot
Rotation orders of the four-dimensional lattice & Crystallographic 1, 2,
3, 4, 6 (§3.4) / 1, 2, 3, 4, 5, 6, 8, 10, 12 of §20.3 of the unified
hypothesis paper {[}S21{]} \\
Decomposition of 137 & 1 + 8 + 128 {[}S19{]} / 1 + 8 + 24 + 40 + 64
{[}S21{]} \\
Oddness & Odd R (G10) / a single state has odd R² (paper 5 of the
dual-geometry series) \\
Number of components per generation & 16 (G15, E4) / 15 {[}S4{]} \\
Order of the exchange & 248 (clock of observation) / 124 (order of the
root) {[}S7{]}. The lock-dynamics paper finds that locking does not
select order 124 {[}S26{]} \\
Parity and statistics & Type of statistics (G14, S17) / a conventional
label \\
Cross term 2ab & Mass (S9) / oscillation of curvature (the
complex-rotating-wave paper) \\
Time reversal and the arrow of time & One side (S12, W1) / the forced
bit is separate from the arrow of time (paper 14 of the dual-geometry
series) \\
Axis Q & Three axes Q₁Q₂Q₃ and 16 states (G15) / one axis with 8 values
and 62 states (W8, the three-layer model) \\
Lorentz signature & Assumed (gnomonic 1) / derived (gnomonic 10, S3) \\
How SU(3) is obtained & Through G₂ (S8) / by the 3 + 2 decomposition
{[}S18{]} \\
\end{longtable}
}

\subsubsection{11.4 Withdrawn or corrected
claims}\label{withdrawn-or-corrected-claims}

When citing, the corrected version is the standard. ``Zero Closure Was
Four-Dimensional'' v4 withdrew ``compactness implies a discrete
spectrum'' {[}S5{]}. Three-layer model 7 and D1 v2.1 withdrew
H\_xyz\_inphase and δδt. Paper 5 of the BH thermodynamics programme
corrected the ``disagreement'' in the abstract of paper 1 as a confusion
between 3 + 1 and 4 + 1. Paper 8 of the BH thermodynamics programme
changed its title to ``structural correspondence'' in v2. Papers 6 and
12 of the dual-geometry series and the observer--system beat conjecture
have also been corrected.

\begin{center}\rule{0.5\linewidth}{0.5pt}\end{center}

\subsection{12. Position among Prior
Studies}\label{position-among-prior-studies}

There are prior studies close to individual elements.

\begin{itemize}
\tightlist
\item
  \textbf{Space from Hilbert space}: Cao, Carroll and Michalakis
  {[}X1{]} assumed a tensor-product decomposition into subsystems and
  assembled space from mutual information. Carroll and Singh {[}X2{]}
  took only the list of eigenvalues of the Hamiltonian and the
  components of one vector as the minimal ontology. Both give time
  evolution from outside and place the tensor-product decomposition
  (labels of subsystems) first; from the side of anonymity, the latter
  is a hidden assumption. Stoica {[}X3{]} criticized that
  three-dimensional space does not emerge uniquely from the minimal
  quantum structure and that a choice is needed. This points to the same
  place as this paper's ``until the projection method is chosen, what is
  readable is undetermined'' (§4.4).
\item
  \textbf{Philosophy of symmetry}: from the classics of Weyl {[}X22{]}
  and Wigner to the collection on the philosophy of symmetry in physics
  {[}X21{]}.
\item
  \textbf{Anonymity}: permutation invariance is considered a postulate
  consistent with quantum theory but not derivable from it {[}X4{]}.
  Structural realism holds that what exists in the world is only the
  structure of relations {[}X5, X6{]}. A Japanese exposition of
  relationalism (Leibniz, Mach), which regards spacetime as relations
  among things, is {[}X33{]}.
\item
  \textbf{Self-consistency}: Chew's bootstrap combined self-consistency
  with the absence of elementary particles (``democracy'') but placed
  spacetime as a background {[}X7{]}.
\item
  \textbf{A world of ratios only}: Barbour's shape dynamics holds that
  only ratios carry objective information and showed that imposing scale
  invariance makes the energy and angular momentum of an island universe
  zero {[}X10, X11{]}. This has the same form as the closure of this
  paper but does not treat discreteness or magic numbers.
\item
  \textbf{Dimensionless numbers}: Duff holds that the number of
  dimensional fundamental constants is zero and that only dimensionless
  numbers are fundamental {[}X13{]}.
\item
  \textbf{Noether's theorem}: the second theorem and strong laws
  {[}X8{]}, quantum anomalies {[}X9{]}, the view of symmetries as
  meta-laws {[}X14{]}.
\item
  \textbf{Duality}: the schema of De Haro and Butterfield {[}X12{]};
  Kaluza--Klein momentum and charge {[}X15{]}.
\item
  \textbf{Triality and three generations}: the proposals of Boyle and of
  Furey and Hughes start from octonions and Clifford algebras and are at
  the stage of conjecture {[}X16, X17{]}.
\item
  \textbf{Attempts to derive α from structure}: Eddington, Rosen, Smith
  and claims from E₈ all lie on the side demanding agreement {[}X18,
  X19, X20{]}.
\item
  \textbf{Locking}: the devil's staircase of the circle map and the
  renormalization of the Farey tree {[}X23, X24{]}; the modular group in
  the quantum Hall effect {[}X25{]}.
\item
  \textbf{Precedents for magic numbers}: nuclear magic numbers
  {[}X30{]}, moonshine {[}X31{]}, e\^{}\{π√163\} {[}X32{]}.
\end{itemize}

Within the range examined, the following combinations were not found:
(1) a construction taking anonymity and self-consistency as axioms and
obtaining closure and return as two theorems; (2) an axiom that what is
attached to a symmetry is not a conserved quantity but a dimensionless
number; (3) an argument combining conjugacy and duality with anonymity
to merge separate symmetries as different views of one symmetry; (4) the
position reading ``near but not equal'' as locking and quasi-oscillation
among several structural numbers; (5) a route into three generations
from the triality q ↦ a q a built into the Hurwitz lattice. Details of
the survey are in Appendix C.

\begin{center}\rule{0.5\linewidth}{0.5pt}\end{center}

\subsection{13. Open Problems}\label{open-problems}

\subsubsection{13.1 The root difficulty}\label{the-root-difficulty}

What is visible depends on the projection method (§4.4). The map from
the state space \{zₙ\} to spacetime \{xyzt\} must be at least a partial
projection. The projection method is unknown.

\subsubsection{13.2 The readout map}\label{the-readout-map}

One zₙ is merely a point in a two-dimensional plane. How do three
independent spatial directions and one time assemble from the relations
among many zₙ? It is not known how a set of complex numbers carries the
quaternion structure required by the condition on composition algebras.
Standard theory proceeded without this map: it accepted discreteness as
a result, gave the background spacetime and the ruler by hand, and
advanced with the computational methods of bra--ket notation, gauge
theory and renormalization. Weyl made R carry the role of Q in 1918 and
discarded R from gauge theory in 1929.

\subsubsection{13.3 The unit of the ruler and the two
orders}\label{the-unit-of-the-ruler-and-the-two-orders}

The closure does not fix the ruler (§3.3). What fixes the unit of R? How
do the order of the exchange and the lattice orders 4 and 6 combine?
These two are different faces of the fact that the readout map is
unknown. Candidates are the conjugacy of amplitude and phase (§4.2) and
the locking of the projection ratio (§4.4).

\subsubsection{13.4 Discipline of numerical
experiments}\label{discipline-of-numerical-experiments}

The catalogue of this paper is also the pass/fail standard for numerical
experiments. The author's model built on the 600-cell imported a
Dirac-type generator, derivatives of position, spin matrices and a
Poisson equation, none of which are in the skeleton, and passed all
stages by swapping parts and readings each time a result did not match.
That pass is not a derivation from the skeleton. Later numerical
experiments are to be built only from the finite operations in the
skeleton (the Clifford translation multiplying a unit quaternion from
the left and the array of component rotations multiplied from the
right); as soon as a part not in the catalogue turns out to be needed,
the work stops and the part is stated explicitly as a new assumption.
The group 2I of the 600-cell contains five-fold and ten-fold rotations
and lies outside the crystallographic restriction, so it does not
satisfy the condition of discreteness of spacetime (§3.4).

\subsubsection{13.5 Checking the consequences of the distinguishability
axiom}\label{checking-the-consequences-of-the-distinguishability-axiom}

The following have not been checked for the distinguishability axiom
(A1′). (1) If the indistinguishability of layers is carried through,
does it require the same definition to hold at the layers below the
square of z′ (z⁴ and beyond)? If so, power sums vanish one after another
and may connect to the return Uⁿ = I. (2) In the exponential
representation E = ln\textbar z\textbar{} + iφ, z ↦ z² is E ↦ 2E, a
discrete scale symmetry with factor 2. It is a semigroup symmetry with
no inverse, outside Noether's theorem. (3) At the squared layer z and −z
cannot be distinguished, and the information of parity (sign reversal at
360°) disappears. The disagreement between the exchange orders 248 and
124 (§11.3) is a candidate for the difference in reading the same
rotation at two layers. (4) The relation between the choice of layer of
the squaring map and the four readings by the combination of t and R
(xyz + t, Q₁Q₂Q₃ + R, xyz + R, Q₁Q₂Q₃ + t) (Appendix D).

\subsubsection{13.6 Checking the correspondences listed as
candidates}\label{checking-the-correspondences-listed-as-candidates}

Whether the groups K1--K7 agree up to renaming alone. Whether the
duality of inside and outside means the exchange of particle and
absorber. Whether triality connects to three generations and to the
relation between bosons and fermions. The correspondence between the 16
sign vectors and the spinor 16 of SO(10). Agreement between the Hopf
winding number and the charge of sign vectors. The interpretations of
numbers in §11.3.

\begin{center}\rule{0.5\linewidth}{0.5pt}\end{center}

\subsection{14. Conclusion}\label{conclusion}

This paper fixed, as a catalogue with statuses and provenances, the
symmetries and magic numbers that a research programme placing no
background spacetime must satisfy. The zeroth axiom was anonymity and
the first axiom self-consistency; the distinguishability axiom was
derived from anonymity; the single definition of closure was the linear
zero closure; the square-sum closure was shown to be the same definition
read at the squared layer; and the return of time evolution was a
theorem. Below them we placed big-picture symmetries, conjugacy and
duality, the dimensionless numbers attached to symmetries and the
projection symmetry, and listed the Standard Model, general relativity,
E₈, the algebra and geometry of this system and the locking symmetries
in the same format.

Mathematically, we confirmed that only real orthogonal maps preserve
both the true square and the norm, and that the 24 Hurwitz units are
cycled 8 + 8 + 8 by the triality q ↦ a q a, which preserves the roots of
E₈. We mapped the 156 published papers to the catalogue, reset the axiom
numbering, revised the earlier claim on the gap of α, and kept the
numerical disagreements as problems admitting several interpretations.

The root difficulty of the map between state space and spacetime lies in
the dependence of what is visible on the projection method. Since
spacetime is stably readable, the map must be at least a partial
projection. In the light of one of the success conditions of the
research project design {[}S14{]}, ``to show that certain requirements
are mathematically incompatible'', this paper lays out each requirement
and its provenance in one table and shows which requirements are already
confirmed, which disagree with each other and which are open. Later
papers and numerical experiments take the numbering and statuses of this
paper as their answer-check standard.

\begin{center}\rule{0.5\linewidth}{0.5pt}\end{center}

\subsection{References}\label{references}

\subsubsection{Self-citations (Concept
DOI)}\label{self-citations-concept-doi}

{[}S1{]} N. Kihara, ``Anonymous Equal-Amplitude Composite-Wave Model:
Basic Axiom System v9,'' 2026-07-24. Concept DOI
10.5281/zenodo.21315735.\\
{[}S2{]} N. Kihara, ``An Axiomatic Framework for Physical Symmetry
Selection by Anonymous Partial Projection: Conditional Identification of
E8 in an Eight-Component Two-Quartet Branch,'' 2026-07-24. Concept DOI
10.5281/zenodo.21521899.\\
{[}S3{]} N. Kihara, ``Symmetry Generation from Zero Closure, Finite
Order, and Self-Consistent Geometry,'' 2026-08-20. Concept DOI
10.5281/zenodo.22028072.\\
{[}S4{]} N. Kihara, ``Noether Conservation Laws and Relational Phase
Dynamics from Discrete Zero Closure,'' 2026-08-21. Concept DOI
10.5281/zenodo.22040735.\\
{[}S5{]} N. Kihara, ``Zero Closure Was Four-Dimensional: An Interpretive
Note on the Geometry and Algebra of Σx\_n² = 0,'' v4, 2026-08-12.
Concept DOI 10.5281/zenodo.21902805.\\
{[}S6{]} N. Kihara, ``The Geometric Identity of the Zero-Square-Sum
Constraint under Scale Invariance,'' 2026-07-22. Concept DOI
10.5281/zenodo.21495305.\\
{[}S7{]} N. Kihara, ``Discovery of Finite-Order Resonance in Iterated
Exchange Scattering,'' 2026-07-18. Concept DOI
10.5281/zenodo.21421366.\\
{[}S8{]} N. Kihara, ``Dual Geometry of Wavelength Space and Frequency
Space (Paper 1),'' 2026-06-08. Concept DOI 10.5281/zenodo.20588036.\\
{[}S9{]} N. Kihara, ``Closed Four-Degree-of-Freedom Structure and Its
Correspondence with 4-Dimensional Lattice Counting (Paper 3),''
2026-06-08. Concept DOI 10.5281/zenodo.20589261.\\
{[}S10{]} N. Kihara, ``Position Had Never Been Read: Sign Sectors,
Duality Theorems, the Readout Hierarchy, and the Derived Clock (Paper
13),'' 2026-06-14. Concept DOI 10.5281/zenodo.20665633.\\
{[}S11{]} N. Kihara, ``Complex Phase Structure Emerging from Nontrivial
Quadratic Closure,'' 2026-07-18. Concept DOI 10.5281/zenodo.21422505.\\
{[}S12{]} N. Kihara, ``Mechanism of Inflation-like Rapid Expansion in
Self-Consistent Closed Relational-Wave Systems,'' second edition,
2026-08-30. Concept DOI 10.5281/zenodo.22112008.\\
{[}S13{]} N. Kihara, ``Anonymity and the Minimal Unit of Information,''
2026-09-12. Concept DOI 10.5281/zenodo.22729127.\\
{[}S14{]} N. Kihara, ``Designing a Research Project to Explore the
Emergence of Spacetime, Particles, and Interactions from States and
Relations,'' 2026-09-20. Concept DOI 10.5281/zenodo.22851944.\\
{[}S15{]} N. Kihara, ``The First Thought Experiment: Tracing the
Equivalence Principle Down to Two Nameless Degrees of Freedom,''
2026-09-20. Concept DOI 10.5281/zenodo.22857952.\\
{[}S16{]} N. Kihara, ``The Second Thought Experiment: Deriving the
Coulomb-Type Inverse-Square Force from a Product Readout and an Area
Clock,'' 2026-09-21. Concept DOI 10.5281/zenodo.22867336.\\
{[}S17{]} N. Kihara, ``The Eleventh Thought Experiment: Emission Records
of Two-Body Systems without a Background Spacetime,'' 2026-10-07.
Concept DOI 10.5281/zenodo.23199299.\\
{[}S18{]} N. Kihara, ``The Twelfth Thought Experiment: From the
Fine-Structure Constant to Observation Maps and Gauge Symmetry,''
2026-10-08. Concept DOI 10.5281/zenodo.23226258.\\
{[}S19{]} N. Kihara, ``The Fourteenth Thought Experiment: Why Can
Different Symmetries Never Be Fully Reconciled?,'' 2026-10-12. Concept
DOI 10.5281/zenodo.23315557.\\
{[}S20{]} N. Kihara, ``Emergence of Discrete Born-Type Weights in
Iterated Two-Channel Exchange Systems,'' 2026-07-18. Concept DOI
10.5281/zenodo.21422470.\\
{[}S21{]} N. Kihara, ``Reading Out Time, Space, Mass, Charge,
Translational Velocity and Spin from the Relations and Interference of
Complex Waves without a Background Spacetime,'' v1.1, 2026-10-15.
Concept DOI 10.5281/zenodo.23384928.\\
{[}S22{]} N. Kihara, ``The Equation of Motion of a Single Wave in a
Relational-Wave Closed System,'' 2026-09-12. Concept DOI
10.5281/zenodo.22728760.\\
{[}S23{]} N. Kihara, ``Analytic Exact Solution of the 90-Degree
High-Symmetry Floor,'' 2026-09-12. Concept DOI
10.5281/zenodo.22728995.\\
{[}S24{]} N. Kihara, ``A Finite Complete-Relation Closed System with
Time Evolution as Vertices: Exact Minimal Closure at N=1, K=4 under
U\^{}K=I,'' 2026-09-15. Concept DOI 10.5281/zenodo.22763329.\\
{[}S25{]} N. Kihara, ``Quadratic Condensation and Relational Local Time
in a Many-Wave Complete Relational System,'' 2026-09-16. Concept DOI
10.5281/zenodo.22788737.\\
{[}S26{]} N. Kihara, ``Lock Dynamics in an Anonymous Two-Channel Closed
Wave System,'' 2026-08-02. Concept DOI 10.5281/zenodo.21763999.\\
{[}S27{]} N. Kihara, ``Selection of the State-Exchange Weight G\_R = 1 −
R and Candidate Correspondence with the Fine-Structure Constant,''
2026-07-16. Concept DOI 10.5281/zenodo.21396760.\\
{[}S28{]} N. Kihara, ``Chain Complex Structure on the 4D Hypercubic
Lattice: Structural Correspondence between Kihara Cube-Packing and
Wilson Lattice Gauge Theory (BH Thermodynamics Programme, Paper 8),''
2026-04-29. Concept DOI 10.5281/zenodo.19880467.\\
{[}S29{]} N. Kihara, ``A Geometric Identity for the Fine-Structure
Constant: From the 4D Unit Ball Volume and its Cube-Packing Deficit (BH
Thermodynamics Programme, Paper 7),'' 2026-04-29. Concept DOI
10.5281/zenodo.19869266.\\
{[}S30{]} N. Kihara, ``BH Thermodynamics Programme Paper 7 Supplement:
Geometric Observations on the Second-Order Correction of the α
Identity,'' 2026-05-01. Concept DOI 10.5281/zenodo.19933729.\\
{[}S31{]} N. Kihara, ``Discrete Structure of a Four-Dimensional Ball:
Unit-Cube Packing and the Asymptotic Volume Deficit Δ(R) (BH
Thermodynamics Programme, Paper 3),'' 2026-04-28. Concept DOI
10.5281/zenodo.19837591.\\
{[}S32{]} N. Kihara, ``Diameter of the Circumscribed Hypersphere of a
Unit Four-Dimensional Hyperrectangle (Gnomonic Projection 10),''
2026-04-28. Concept DOI 10.5281/zenodo.19533312.

\subsubsection{External literature}\label{external-literature}

{[}X1{]} C. Cao, S. M. Carroll, S. Michalakis, ``Space from Hilbert
Space: Recovering Geometry from Bulk Entanglement,'' Phys. Rev.~D 95
(2017); arXiv:1606.08444.\\
{[}X2{]} S. M. Carroll, A. Singh, ``Mad-Dog Everettianism: Quantum
Mechanics at Its Most Minimal,'' arXiv:1801.08132 (2018).\\
{[}X3{]} O. C. Stoica, ``3D-Space and the preferred basis cannot
uniquely emerge from the quantum structure,'' arXiv:2102.08620.\\
{[}X4{]} ``Quantum Statistics of Identical Particles,'' Foundations of
Physics 52:77 (2022).\\
{[}X5{]} S. French, J. Ladyman, ``Remodelling structural realism,''
Synthese (2003).\\
{[}X6{]} J. Ladyman, D. Ross, \emph{Every Thing Must Go} (Oxford
University Press, 2007).\\
{[}X7{]} G. F. Chew, the bootstrap programme (1959--); ``Bootstrap
model,'' Wikipedia.\\
{[}X8{]} Y. Kosmann-Schwarzbach, \emph{The Noether Theorems} (Springer,
2011); ``The Noether theorems in context,'' arXiv:2004.09254.\\
{[}X9{]} ``Anomalies,'' in \emph{Renormalization} (Cambridge University
Press).\\
{[}X10{]} J. Barbour, ``Scale-invariant gravity: particle dynamics,''
arXiv:gr-qc/0211021.\\
{[}X11{]} J. Barbour, ``Shape dynamics. An introduction,''
arXiv:1105.0183; ``The Physics and Metaphysics of Pure Shape Dynamics,''
arXiv:2203.02773.\\
{[}X12{]} S. De Haro, J. Butterfield, ``A Schema for Duality,
Illustrated by Bosonization,'' arXiv:1707.06681; \emph{The Philosophy
and Physics of Duality} (OUP, 2025).\\
{[}X13{]} M. J. Duff, L. B. Okun, G. Veneziano, ``Trialogue on the
number of fundamental constants,'' JHEP (2002); arXiv:physics/0110060.\\
{[}X14{]} M. Lange, ``Are symmetry principles meta-laws?'' Synthese
(2025).\\
{[}X15{]} ``A New Perspective on Kaluza-Klein Theories,''
arXiv:2404.05302.\\
{[}X16{]} L. Boyle et al., ``The C\(\ell\)(8) algebra of three fermion
generations with spin and full internal symmetries,''
arXiv:1906.05102.\\
{[}X17{]} ``Sedenions, Clifford Algebras, and Three Fermion Generations:
A Focused Review,'' arXiv:2609.10569; ``The Standard Model, The
Exceptional Jordan Algebra, and Triality,'' arXiv:2006.16265.\\
{[}X18{]} H. Kragh, ``Magic Number: A Partial History of the
Fine-Structure Constant,'' Archive for History of Exact Sciences 57
(2003) 395--431.\\
{[}X19{]} U. D. Jentschura, I. Nandori, ``Attempts at a determination of
the fine-structure constant from first principles: A brief historical
overview,'' Eur. Phys. J. H 39 (2014) 591--613; arXiv:1411.4673.\\
{[}X20{]} ``On Arthur Eddington's Theory of Everything,''
arXiv:1510.04046.\\
{[}X21{]} K. Brading, E. Castellani (eds.), \emph{Symmetries in Physics:
Philosophical Reflections} (Cambridge University Press, 2003).\\
{[}X22{]} H. Weyl, \emph{Symmetry} (Princeton University Press, 1952).\\
{[}X23{]} P. Cvitanović, M. H. Jensen, L. P. Kadanoff, I. Procaccia,
``Renormalization, unstable manifolds, and the fractal structure of mode
locking,'' Phys. Rev.~Lett. (1985).\\
{[}X24{]} P. Cvitanović, B. Shraiman, B. Söderberg, ``Scaling laws for
mode lockings in circle maps,'' Physica Scripta (1985).\\
{[}X25{]} B. P. Dolan, ``Duality and the Modular Group in the Quantum
Hall Effect,'' J. Phys. A 32 (1999) L243; arXiv:cond-mat/9805171.\\
{[}X26{]} V. Elser, N. J. A. Sloane, ``A highly symmetric
four-dimensional quasicrystal,'' J. Phys. A (1987).\\
{[}X27{]} J. A. Wheeler, R. P. Feynman, ``Interaction with the Absorber
as the Mechanism of Radiation,'' Rev.~Mod. Phys. 17 (1945); ``Classical
Electrodynamics in Terms of Direct Interparticle Action,'' Rev.~Mod.
Phys. 21 (1949).\\
{[}X28{]} J. C. Baez, ``The Octonions,'' Bull. Amer. Math. Soc. 39
(2002).\\
{[}X29{]} J. H. Conway, N. J. A. Sloane, \emph{Sphere Packings, Lattices
and Groups} (Springer).\\
{[}X30{]} D. Bonatsos et al., ``Magic numbers of cylindrical symmetry,''
arXiv:1909.00233; ``Magic numbers for shape coexistence,''
arXiv:1811.01071.\\
{[}X31{]} ``Monstrous moonshine and the classification of CFT,''
arXiv:math/9906167; ``Introduction to Sporadic Groups for physicists,''
arXiv:1305.5974.\\
{[}X32{]} ``Lectures on modular forms and strings,'' arXiv:2208.07242.\\
{[}X33{]} S. Uchii, \emph{Kūkan no nazo, jikan no nazo} {[}Mysteries of
space, mysteries of time{]} (Chūkō Shinsho, 2006), in Japanese.

\begin{center}\rule{0.5\linewidth}{0.5pt}\end{center}

\subsection{Appendix A: Reproducing the
Checks}\label{appendix-a-reproducing-the-checks}

The programs and results of Checks 1, 2 and 2b are in the repository
folder
\texttt{次元の生成構造/\allowbreak{}四次元球殻住民\_\allowbreak{}自己無撞着\_\allowbreak{}視野角三角形\_\allowbreak{}20261016/\allowbreak{}見落とされた対称性の確認\_\allowbreak{}20261017/\allowbreak{}}.
They are reproduced by \texttt{./\allowbreak{}run\_\allowbreak{}all.sh}, and the results are
written as JSON to \texttt{results/\allowbreak{}}.

\subsection{Appendix B: Self-Citation Correspondence
Table}\label{appendix-b-self-citation-correspondence-table}

The correspondence between the 156 published papers and the catalogue
rows, the disagreements and the correspondence with the old numbering
are in \texttt{自己引用の対応表\_\allowbreak{}20261017.md} (self-citation
correspondence table, in the same folder as this paper; in Japanese).

\subsection{Appendix C: The Original Lists and the Prior-Study
Survey}\label{appendix-c-the-original-lists-and-the-prior-study-survey}

The working documents on which this paper is based are (in Japanese):

\begin{itemize}
\tightlist
\item
  \texttt{次元の生成構造/\allowbreak{}対称性の一覧\_\allowbreak{}保存量と固定点\_\allowbreak{}20261017.md}
  (list of symmetries)
\item
  \texttt{次元の生成構造/\allowbreak{}マジックナンバーの一覧\_\allowbreak{}20261017.md} (list of
  magic numbers)
\item
  \texttt{次元の生成構造/\allowbreak{}対称性の一覧\_\allowbreak{}先行研究調査\_\allowbreak{}20261017.md}
  (survey of prior studies)
\end{itemize}

\subsection{Appendix D: Dialogue Record of the Thought Experiment on
Linear Zero
Closure}\label{appendix-d-dialogue-record-of-the-thought-experiment-on-linear-zero-closure}

The distinguishability axiom and the definition of linear closure in
§3.0 came from the dialogue between the author and Claude on 2026-10-18.
The dialogue record is in
\texttt{次元の生成構造/\allowbreak{}線形和のゼロ閉塞\_\allowbreak{}思考実験\_\allowbreak{}20261018/\allowbreak{}線形和のゼロ閉塞Σzn=0仮定の追加の思考実験\_\allowbreak{}20261018.md}
(no summarization; only typos in the author's remarks corrected; in
Japanese).

\subsection{Version History}\label{version-history}

\begin{itemize}
\tightlist
\item
  v1.0 (2026-10-18): first edition. Version DOI 10.5281/zenodo.23444452,
  Concept DOI 10.5281/zenodo.23444451.
\end{itemize}

\end{document}
